\subsection{Магазин}

\subsubsection{Описание решения}

Воспользуемся жадной стратегией: максимизируем стоимость самых дешёвых позиций в каждом чеке. Для определённости стоит отсортировать данный массив по возрастанию.

Если размер чека будет больше \( k \), то массив получится разделить на меньшее число чеков без потери скидки. Ещё при таком подходе большой чек получит более дешёвые позиции, что уменьшает суммарную скидку. По этой причине не имеет смысла делить массив на части, кратные \( k \).

В итоге остаётся делить массив на чеки размера \( k \) \textit{(кроме, возможно, последнего)}, начиная с правой части. В скидку пойдут крайние левые позиции из каждого чека размера \( k \) --- она и будет максимальной.

\subsubsection{Асимптотическая сложность}

\textit{Временная сложность решения} --- линейно-логарифмическая \( \mathcal{O}(N\cdot\log N) \). В алгоритме используется реализация сортировки из \textit{STL}, которая имеет линейно-логарифмическую сложность согласно официальной документации.

\textit{Пространственная сложность решения} --- линейная \( \mathcal{O}(N) \). Для сортировки массива требуется сначала его сохранить в памяти.

\subsubsection{Листинг кода}

\begin{lstlisting}
#include <algorithm>
#include <iostream>
#include <vector>

int main() {
  int n, k;
  std::cin >> n >> k;
  std::vector<int> arr;
  arr.reserve(n);
  long sum = 0;
  for (int i = 0; i < n; i++) {
    int item;
    std::cin >> item;
    arr.push_back(item);
    sum += item;
  }
  std::sort(arr.begin(), arr.end());
  for (int i = n - k; i >= 0; i -= k) {
    sum -= arr[i];
  }
  std::cout << sum;
  return 0;
}
\end{lstlisting}